% 组合数学（综述）
% license CCBYSA3
% type Wiki

本文根据 CC-BY-SA 协议转载翻译自维基百科\href{https://en.wikipedia.org/wiki/Combinatorics}{相关文章}。

组合数学是数学的一个分支，主要研究计数（counting）——既作为获得结果的方法，又作为研究有限结构某些性质的目的。它与数学的许多其他领域密切相关，并在从逻辑到统计物理、从进化生物学到计算机科学等众多领域中有广泛的应用。

组合数学以其所处理问题的广泛性而著称。组合问题出现在纯数学的许多分支中，尤其是在代数学、概率论、拓扑学与几何学中\(^\text{[1]}\) ，同时也出现在众多应用领域。  
历史上，许多组合问题往往是独立提出的，作为在某个数学背景下出现的特定问题的特设性解法（ad hoc solutions）。然而，在二十世纪后期，强大而一般的理论方法被发展出来，使得组合数学成为一门独立的数学分支\(^\text{[2]}\)。组合数学中最古老且最易于理解的部分之一是图论（graph theory），它本身就与其他数学领域有着大量自然的联系。在计算机科学中，组合数学被广泛用于推导算法分析中的公式与估计。
\subsection{定义}
组合数学的确切范围并没有被普遍认同\(^\text{[3]}\)。正如 H. J. Ryser 所指出的，由于组合数学跨越了众多数学分支，因此要给出一个精确的定义是困难的\(^\text{[4]}\)。如果我们从其所研究的问题类型来描述这一领域，那么组合数学主要涉及以下几个方面：
\begin{enumerate}
  \item 枚举（enumeration）：对特定结构的计数，这些结构在广义上被称为排列（arrangements）或配置（configurations），它们通常与有限系统相关；
  \item 存在性（existence）：研究满足特定条件的结构是否存在；
  \item 构造性（construction）：探讨如何以一种或多种方式构造出这些结构；
  \item 优化（optimization）：在多种可能的结构或解中寻找“最优”的一个，无论其意义为“最大”、“最小”或符合其他最优性标准。
\end{enumerate}
Leon Mirsky 曾指出：“组合数学是一系列彼此关联的研究，它们有共同点，但在目标、方法以及所达到的统一程度上却差异很大。”\(^\text{[5]}\)一种定义组合数学的方式，或许是通过描述它的各个分支（subdivisions）、它们所处理的问题以及所采用的技术。下文将采用这种方式来介绍组合数学。然而，从历史的角度看，是否将某些主题纳入组合数学的范畴，也存在一定的传统性和争议性\(^\text{[6]}\)。尽管组合数学主要关注有限系统，但其中的一些问题与方法也可以推广到无限（尤其是可数）离散系统的情形。

